
\subsection{非消失定理与马斯洛夫类猜想}
\label{subsec:6-1-2}

回忆 $L$ 的基本循环的庞加莱对偶
是我们滤过 $A_\infty$ 代数的同伦幺元。
利用这一点，
可以证明弗洛尔上同调的非消失定理。
本小节也讨论马斯洛夫类猜想\liftIndex{I022}。

\begin{theorem}{6.1.9}\label{thm:6-1-9}
设 $L\subset M$ 是相对 spin 拉格朗日子流形，
而 $\mathbf b=(\boldsymbol b,b)
\in\mathcal M_{\mathrm{weak,def}}(L)$\translatornote{原书实体页 366
在这里将下标写成 $\mathrm{def,weak}$，
与本书其他各处的 $\mathrm{weak,def}$ 次序不一致；
此处统一为后者。见疑误 A-0356。}。
则定理 D 的谱序列在
$\mathbf b_1=\mathbf b_0=\mathbf b$ 时
具有下列性质。

用
$PD([L])\in H^0(L;\mathbb Q)$、
$PD([pt])\in H^n(L;\mathbb Q)$
分别表示 $L$ 的基本循环和点类的庞加莱对偶。
将它们视为
$E_2=H(L;\mathbb Q)\otimes
\operatorname{gr}_*(\mathfrak F\Lambda_{0,\mathrm{nov}})$
中的项。则有：

\par\smallskip\noindent\textup{(6.1.10)}\quad
对每个 $r$，
\begin{align}
\delta_r\bigl(
 PD([L])-\mathfrak h_2^{+M}
 (e^{\boldsymbol b},e^b\mathbin{\overline\otimes}e^b)
 \bigr)&=0,
\tag{6.1.10.1}\label{eq:6-1-10-1}\\
PD([pt])&\notin\operatorname{Im}(\delta_r).
\tag{6.1.10.2}\label{eq:6-1-10-2}
\end{align}
\textup{(6.1.10)} 对
$\Lambda_{\mathrm{nov}}$ 系数也成立。

\pagebreak[3]
进一步作如下假设。
\begingroup\upshape
\begin{statement}{假设}{6.1.11}\label{ass:6-1-11}
凡边界落在 $L$ 上的伪全纯圆盘，
其马斯洛夫指标均非正。
\end{statement}
\endgroup

\par\smallskip\noindent\textup{(6.1.12)}\quad
于是，对每个 $r$，
\begin{align}
PD([L])-\mathfrak h_2^{+M}
 (e^{\boldsymbol b},e^b\mathbin{\overline\otimes}e^b)
 &\notin\operatorname{Im}(\delta_r),
\tag{6.1.12.1}\label{eq:6-1-12-1}\\
\delta_r(PD([pt]))&=0.
\tag{6.1.12.2}\label{eq:6-1-12-2}
\end{align}
\textup{(6.1.12)} 对
$\Lambda_{\mathrm{nov}}$ 系数也成立。
因此，在假设 6.1.11 下，
$PD([pt])$ 与 $PD([L])$
确定 $E_\infty$ 中线性无关的元素。
特别地，
\[
HF((L,\mathbf b);\Lambda_{\mathrm{nov}})
\ne\{0\}.
\]
\end{theorem}

$\mathfrak h_2^{+M}$ 的记号见
第 3.8.4 小节的 (3.8.35.4)。
注意，在假设 6.1.11 下，
bulk 形变后的弱无障碍性
等价于 bulk 形变后的无障碍性。
假设 6.1.11 还可以由下列拓扑条件代替。

\begin{statement}{假设}{6.1.13}\label{ass:6-1-13}
若 $\beta\in\pi_2(M,L)$ 且
$\mu_L(\beta)>0$，
则 $\omega(\beta)<0$。
\end{statement}

导论中的定理 E 是定理 6.1.9 的直接推论。
当 $\mathbf b_1\ne\mathbf b_0$ 时，
定理 6.1.9 有如下版本。

\begin{theorem}{6.1.14}\label{thm:6-1-14}
设 $\mathbf b_1,\mathbf b_0
\in\mathcal M_{\mathrm{weak,def}}(L)$，
满足
\[
\mathfrak{PO}(\mathbf b_1)=\mathfrak{PO}(\mathbf b_0),
\qquad
\pi_{\mathrm{amb}}(\mathbf b_1)
=\pi_{\mathrm{amb}}(\mathbf b_0),
\]
其中 $\mathbf b_i=(\boldsymbol b,b_i)$。
则
\begin{equation}
\delta_r\bigl(
PD([L])-\mathfrak h_2^{+M}
(e^{\boldsymbol b},e^{b_1}\mathbin{\overline\otimes}e^{b_0})
\bigr)=0.
\tag{6.1.10.1$'$}\label{eq:6-1-10-1-prime}
\end{equation}
若再加上假设 6.1.11，
则
\begin{align}
PD([L])-\mathfrak h_2^{+M}
(e^{\boldsymbol b},e^{b_1}\mathbin{\overline\otimes}e^{b_0})
&\notin\operatorname{Im}(\delta_r),
\tag{6.1.12.1$'$}\label{eq:6-1-12-1-prime}\\
\delta_r(PD([pt]))&=0.
\tag{6.1.12.2$'$}\label{eq:6-1-12-2-prime}
\end{align}
\textup{(6.1.12.1$'$)} 对
$\Lambda_{\mathrm{nov}}$ 系数也成立。
\end{theorem}

作为非消失定理的应用，
可得下列两个定理；
它们将在第 6.4.4 小节证明。

\begin{theorem}{6.1.15}\label{thm:6-1-15}
设 $L\subset M$ 是相对 spin 拉格朗日子流形。
若限制映射
$H^2(M;\mathbb Q)\to H^2(L;\mathbb Q)$
满射，且 $L$ 的马斯洛夫指标
在 $\pi_2(M,L)$ 上恒为零，
则 $L$ 在 bulk 形变后无障碍，
并且对任意
$\mathbf b\in\mathcal M_{\mathrm{def}}(L)$，
都有
\[
HF((L,\mathbf b),(L,\mathbf b);
\Lambda_{\mathrm{nov}})\ne0.
\]
\end{theorem}

注意，马斯洛夫指标恒为零的假设
蕴含 $c_1(M)$ 在 $\pi_2(M)$ 上为零。
例如，特殊拉格朗日子流形的
马斯洛夫指标为零 \textbf{[Daz81]}。
把本定理与托马斯、丘成桐的结果
\textbf{[ThYa02]} 结合可知：
若 $L$ 满足定理 6.1.15 的条件，
则其哈密顿形变类中
至多有一个特殊拉格朗日子流形。

由于 $\Lambda_{\mathrm{nov}}$ 系数的
弗洛尔上同调在哈密顿形变下不变，
立刻得到关于所谓马斯洛夫指标猜想的推论。

\begin{corollary}{6.1.16}\label{cor:6-1-16}
设 $L\subset M$ 是相对 spin 拉格朗日子流形，
$H^2(M;\mathbb Q)\to H^2(L;\mathbb Q)$
满射，且存在哈密顿微分同胚 $\psi$，
使 $L\cap\psi(L)=\varnothing$。
则 $L$ 的马斯洛夫指标并非恒为零。
\end{corollary}

还可以估计极小马斯洛夫数。
记 $\Sigma_L$ 为 $L$ 的极小马斯洛夫数；
它是由
\[
\operatorname{Im}(\mu_L)=\Sigma_L\mathbb Z,
\qquad
\mu_L:\pi_2(M,L)\longrightarrow\mathbb Z
\]
定义的非负整数。

\begin{theorem}{6.1.17}\label{thm:6-1-17}
设 $L\subset M$ 是 $n$ 维相对 spin
拉格朗日子流形，
存在哈密顿微分同胚 $\psi$，
使 $L\cap\psi(L)=\varnothing$。
\begin{enumerate}
\item[(1)] 若 $L$ 在 bulk 形变后弱无障碍，
则 $\Sigma_L\ne0$。
\item[(2)] 若还存在
$(\boldsymbol b,b)\in
\widehat{\mathcal M}_{\mathrm{weak,def}}(L)$，
满足
\[
\boldsymbol b\in H^2(M;\mathbb C)
\otimes\Lambda_{0,\mathrm{nov}}^+,
\qquad
b\in H^1(L;\mathbb C)
\otimes\Lambda_{0,\mathrm{nov}}^+,
\]
则 $2\le\Sigma_L\le n+1$。
\end{enumerate}
\end{theorem}

这里的 $M$ 不必闭。
只要 $L$ 紧，
本书的结果也适用于格罗莫夫意义下的
驯服 $M$ \textbf{[Grom85]}，
例如 $M=\mathbb C^n$。
推论 6.1.16 因而蕴含导论中的定理 K。
当 $M$ 球面正时，
极小马斯洛夫数还将在
\textbf{[FOOO09I]} 中进一步讨论。

现在回顾紧拉格朗日嵌入
$L\subset\mathbb C^n$
的马斯洛夫类相关结果的历史。
在 $\mathbb C^n$ 中，
同态
$\mu_L:\pi_2(\mathbb C^n,L)\to\mathbb Z$
下降为一个上同调类
$\mu_L\in H^1(L;\mathbb Z)$；
这就是阿诺尔德 \textbf{[Arn67]}
定义的标准马斯洛夫类。

关于 $\mathbb C^n$ 中 $L$ 的马斯洛夫类
$\mu_L$，有两个重要猜想。
第一个通常称为马斯洛夫类猜想：
任意紧拉格朗日嵌入
$L\subset\mathbb C^n$
的马斯洛夫类 $\mu_L$ 都非零。
此前已知的只有环面
$L\cong T^n$，
或者更一般地，
$L\subset\mathbb C^n$
具有非正曲率度量的情形，
由克洛德·维泰博（Claude Viterbo）
\textbf{[Vit90]} 证明；
以及 $\mathbb C^2$ 中的二维情形
\textbf{[Pol91I]}。
这个猜想还可加强为：
对任意可位移的拉格朗日子流形\liftIndex{I219}
$L\subset M$，
存在 $\beta\in\pi_2(M,L)$，
使 $E(\beta)>0$ 且 $\mu_L(\beta)>0$。
这里“可位移”指存在哈密顿微分同胚
$\phi:M\to M$，
使 $\phi(L)\cap L=\varnothing$。
事实上，\textbf{[Vit90]} 与
\textbf{[Pol91I]}
证明的都是这一较强断言。

第二个猜想或问题通常称为
奥丹问题\liftIndex{I216}
\textbf{[Aud88]}：
在 $\mathbb C^n$ 中，
任一嵌入的拉格朗日环面
都有 $\Sigma_L=2$。
当 $n=2$ 时，
这也由 \textbf{[Vit90, Pol91I]} 得出。
第二作者吴龙根在
\textbf{[Oh96I]} 中证明，
对 $n\le24$，
单调拉格朗日环面
$T^n\subset\mathbb C^n$
满足此猜想。

本书的 2000 年版本
流传以后，
又出现或公布了若干结果。
首先，我们将利用谱序列与乘积结构的相容性
\textup{(D.4)}，
在定理 6.4.35 中
去掉前述维数限制，
从而对单调拉格朗日环面
完全解答奥丹问题。
使用这一相容性的想法来自塞德尔
\textbf{[Sei01]}，
另见 \textbf{[Bi04, Buh06]}。
我们将在第 6.4 节末讨论这一点。

借助辛场论 \textbf{[EGH00]} 的想法，
雅科夫·埃利亚什伯格（Yakov Eliashberg）
提出了肯定回答奥丹问题的方案
\textbf{[Eli01]}；
齐利巴克与 K. Mohnke
\textbf{[CiMo04]}
进一步详述了该方案。
第一作者深谷贤治
\textbf{[Fuk05II]} 还概述了另一方案：
结合弦拓扑 \textbf{[ChSu99]} 的想法，
拉格朗日弗洛尔理论也能用于证明
奥丹问题及有关
aspherical 拉格朗日子流形的其他结果。
